\documentclass{article}
\usepackage{iclr2027_conference,times}
\input{math_commands.tex}
\usepackage{booktabs}
\usepackage{graphicx}
\usepackage{hyperref}
\usepackage{url}

\title{Tensor-Core Fast Walsh--Hadamard Transforms\\at the Memory Roofline}

% Keep authors anonymous for submission. Do not enable \iclrfinalcopy.
\author{Anonymous authors}

\begin{document}
\maketitle

\begin{abstract}
The Fast Walsh--Hadamard Transform (FWHT) is increasingly used to suppress
outliers before low-precision quantization, but conventional GPU kernels pay
substantial instruction and data-movement overhead. We present a tensor-core
FWHT whose register layout and transform ordering reduce dynamic instructions
while preserving coalesced movement through the memory hierarchy. On an
NVIDIA H100, the implementation sustains approximately 3.0~TB/s effective
bandwidth, operates within 1.7\% of an empirical read--modify--write roofline,
and is up to 1.59$\times$, 2.36$\times$, 2.25$\times$, and 1.36$\times$ faster
than HadaCore, Tri Dao's CUDA implementation, FlagGems Triton, and a patched
Triton tensor-core baseline, respectively. Nsight Compute shows
2.02--3.60$\times$ fewer dynamically executed warp instructions than HadaCore
across representative transform sizes. We also fuse a randomized H16 transform
with Transformer Engine block-FP8 quantization. On an H100 NVL and a
Llama-shaped MLP, the
fused path is 1.974$\times$ faster for forward plus backward than a
functionality-matched composition of unmodified Tri Dao H16, stock PyTorch
signs, and unmodified Transformer Engine, while adding 5.7\% latency relative
to Transformer Engine with the RHT omitted. Native NVFP4 results remain pending.
\end{abstract}

\section{Introduction}
Hadamard rotations are a simple and effective way to redistribute outliers
before aggressive quantization \citep{TODO_hadamard_quantization}. This makes
their systems cost important:
an otherwise inexpensive rotation may require extra reads and writes around a
much faster FP8 or FP4 matrix multiplication. We ask whether tensor cores can
execute a standalone FWHT at the memory roofline and whether the same layout
can be integrated into a randomized Hadamard transform (RHT) pipeline.

Our current contributions are:
\begin{itemize}
  \item A tensor-core FWHT supporting normalized FP16 and BF16 transforms of
        size 256 through 32768.
  \item A matched H100 evaluation against four open-source CUDA and Triton
        implementations, including correctness and empirical roofline checks.
  \item Dynamic profiling evidence that the transform ordering executes
        2.02--3.60$\times$ fewer warp instructions than HadaCore.
  \item A fused H16 RHT--FP8 Transformer Engine adapter that is 1.974$\times$
        faster than an off-the-shelf RHT-enabled composition on a Llama-shaped
        MLP, with an explicit no-RHT control.
\end{itemize}

\section{Background and motivation}
For a power-of-two dimension $d$, the normalized Hadamard matrix satisfies
$H_d H_d^\top / d = I$. A randomized Hadamard transform introduces a fixed
Rademacher diagonal matrix $S_d$,
\begin{equation}
  R_d = \frac{1}{\sqrt d} S_d H_d,
\end{equation}
which smooths tensor outliers while preserving inner products. Low-precision
training systems can apply compatible rotations to both GEMM operands so that
the rotations cancel in exact arithmetic. NVIDIA's NVFP4 recipe applies a
tiled $d=16$ RHT to BF16 weight-gradient inputs before quantization
\citep{nvidia_nvfp4}.

Hopper supports native tensor-wide FP8, whereas native MXFP8 and NVFP4 require
Blackwell-class hardware. We therefore use Hopper to evaluate the FWHT,
profiling, fused RHT mechanics, and an FP8 quantization pipeline; native NVFP4
end-to-end measurements are reserved for Blackwell.

\section{Related work}
GPU FWHT implementations span CUDA and Triton kernels, including HadaCore,
Tri Dao's CUDA implementation, FlagGems, and tensor-core decompositions
\citep{TODO_hadacore,TODO_tri_dao_fwht,TODO_flaggems,TODO_arthur_fwht}.
Hadamard rotations are also used to reduce quantization outliers in low-precision
inference and training \citep{TODO_hadamard_quantization}. Transformer Engine
provides the FP8 and NVFP4 training paths used in our application study
\citep{TODO_transformer_engine_fp8}. \utkarsh{TODO: Complete related work and
separate standalone FWHT kernels from rotation-aware quantization systems.}

\section{Tensor-core FWHT}
\label{sec:method}
The implementation decomposes the transform into tensor-core-compatible
subtransforms and selects a register permutation that makes the intermediate
layout useful to the next stage. Small transforms remain entirely in
registers. Larger transforms use shared memory between two phases and overlap
global-to-shared movement with computation. Normalization is folded into the
constant Hadamard operands.

\utkarsh{TODO: Add the register-layout and transform-order diagram; move full
tiling, permutation, and pipeline details to the appendix.}

The central systems hypothesis is that layout selection removes explicit
permutation, address-generation, and data-movement instructions. Section
\ref{sec:profile} tests this hypothesis with dynamic hardware counters.

\section{Evaluation}
\label{sec:evaluation}
\paragraph{Setup.}
We evaluate an NVIDIA H100 80GB HBM3 GPU with PyTorch 2.9.1+cu128 and CUDA
12.9.1 for extension compilation. Each saturated-throughput point processes
$2^{30}$ FP16 or BF16 elements for 400 timed in-place transforms. All baselines
use the same normalization, shape, stream semantics, and preallocated tensors.
Compilation and JIT time are excluded.

\utkarsh{TODO: Define effective bandwidth and the empirical AddOne roofline,
report correctness tolerances, and repeat the standalone table across
independent processes with uncertainty, clocks, power, driver, and compiler
flags.}

\begin{table}[t]
\caption{Speedup of our kernel over matched baselines on H100. Values below
1 indicate a baseline win.}
\label{tab:baselines}
\centering
\small
\begin{tabular}{llrrrr}
\toprule
Type & Size & HadaCore & Tri Dao & FlagGems & Arthur Triton \\
\midrule
FP16 & 256   & 0.99 & 1.77 & 1.02 & 1.81 \\
FP16 & 4096  & 0.99 & 1.28 & 1.15 & 1.11 \\
FP16 & 16384 & 1.10 & 1.33 & 1.52 & 0.98 \\
FP16 & 32768 & 1.53 & 2.36 & 2.23 & 1.28 \\
BF16 & 256   & 0.99 & 1.77 & 1.02 & 1.80 \\
BF16 & 4096  & 1.01 & 1.26 & 1.14 & 1.16 \\
BF16 & 16384 & 1.15 & 1.40 & 1.53 & 0.99 \\
BF16 & 32768 & 1.59 & 2.15 & 2.25 & 1.36 \\
\bottomrule
\end{tabular}
\end{table}

Our kernel sustains 2.99--3.03~TB/s effective bandwidth and remains within
0--1.7\% of an empirical elementwise read--modify--write roofline across the
tested domain. It is not uniformly faster: several small and middle-size
points are ties or small losses. The advantage is largest at size 32768, where
instruction and intermediate-layout costs dominate the baselines.

\subsection{Instruction-level explanation}
\label{sec:profile}
We profile one NVTX-marked launch after warmup with Nsight Compute 2025.2.1.
Table~\ref{tab:profile} reports dynamic warp instructions rather than static
SASS line counts.

\begin{table}[t]
\caption{Matched H100 profile against HadaCore. Instruction reduction and
kernel speedup are HadaCore divided by ours.}
\label{tab:profile}
\centering
\small
\begin{tabular}{llrr}
\toprule
Type & Size & Fewer instructions & Kernel speedup \\
\midrule
FP16 & 256   & 2.92$\times$ & 1.01$\times$ \\
FP16 & 4096  & 2.25$\times$ & 1.04$\times$ \\
FP16 & 32768 & 3.47$\times$ & 1.69$\times$ \\
BF16 & 256   & 2.79$\times$ & 1.01$\times$ \\
BF16 & 4096  & 2.02$\times$ & 1.08$\times$ \\
BF16 & 32768 & 3.60$\times$ & 1.81$\times$ \\
\bottomrule
\end{tabular}
\end{table}

At size 256 the instruction reduction does not improve latency because memory
traffic and launch overhead dominate. At size 32768, FP16 tensor-pipe activity
increases from 9.61\% to 16.54\%, while the profiled kernel becomes
1.69$\times$ faster. BF16 shows the same trend.

\section{Randomized Hadamard integration}
\label{sec:rht}
Our target pipeline is
\begin{equation}
\text{BF16 input}\rightarrow\text{random sign}\rightarrow
\text{Hadamard}\rightarrow\operatorname{amax}\rightarrow
\text{scale}\rightarrow\text{low-precision output}.
\end{equation}
The random sign operation is an exact sign-bit change for BF16/FP16 and can be
folded into loading rather than materialized as a separate tensor. The
normalization is already folded into the transform. The remaining opportunity
is to combine output storage with amax/scaling and quantization, eliminating
full-tensor intermediate memory passes. Section~\ref{sec:rht-results} evaluates
this design.

\paragraph{Transformer Engine adapter.}
One 128-by-128 Triton tile applies eight H16 transforms per block and emits
Transformer Engine's independently scaled rowwise and columnwise block-FP8
views. The backward path applies the exact transpose $R^\top=HS/4$ to Dgrad;
the SwiGLU variant fuses this inverse transform with dSwiGLU.

\paragraph{Dynamic weights and optimizer.}
FP32 AdamW masters and moments remain in the original basis. Each step writes
$WR$ directly into Transformer Engine's block-FP8 weight storage and maps BF16
Wgrad through $R^\top$. The final optimizer kernel combines inverse H16, FP32
moment updates, weight decay, and the original-basis parameter update.

\section{RHT--FP8 evaluation}
\label{sec:rht-results}

\utkarsh{TODO: Add the H100 NVL setup: PyTorch, CUDA, Transformer Engine,
dtype, exact MLP shapes, warmups, iterations, timing boundaries, and profiler
configuration.}

The H100 fusion prototype processes 67.1M elements. Plain H16 RHT is
within 2.3\% of a dense H16 matmul reference and matches it exactly at the
output dtype. Fusing RHT, per-16 amax/scaling, and E4M3 output reduces latency
from 1.578 to 0.262~ms for FP16 input and from 1.580 to 0.264~ms for BF16 input,
a 6.01$\times$/5.98$\times$ speedup over an eager separate-operation pipeline.
This is a fusion ablation rather than a Transformer Engine comparison, and
E4M3 on H100 is not native NVFP4.

Separately, we benchmark Transformer Engine 2.18 linear layers with a matched
PyTorch 2.8.0+cu129 stack on the same H100. For a
$4096\times8192\times8192$ GEMM, delayed-scaling FP8 is 1.32$\times$ faster
than BF16 in forward and 1.66$\times$ faster for the
combined forward, Dgrad, and Wgrad measurement. For Llama-style
$4096\rightarrow11008$ and $11008\rightarrow4096$ projections, the combined
speedups are 1.47$\times$ and 1.41$\times$. FP8 block scaling reaches
1.17$\times$/1.34$\times$ for forward/combined measurement at the square
8192 shape and 1.26$\times$ combined for both Llama projections. Conversely,
all FP8 recipes trail BF16 for the smaller $4096^3$ forward GEMM, showing that
quantization overhead creates a workload-size crossover. These linear-layer
microbenchmarks establish component-level controls for the directly integrated
path below; they are not end-to-end model results.

A forward-only integration-cost experiment applies the fixed-sign H16 transform
to BF16 activations immediately before the same Transformer Engine Linear. The
delayed-scaling FP8 pipeline, including RHT, is 1.19$\times$ faster than plain
BF16 Linear for square-8192 and 1.13$\times$ faster for the Llama up
projection. It is approximately tied for Llama down (0.98$\times$) and slower
for square-4096. When both paths include RHT, delayed FP8 is 1.30$\times$,
1.27$\times$, and 1.13$\times$ faster for square-8192, Llama up, and Llama
down. This forward-only measurement is retained as a component ablation; the
directly integrated forward and backward path follows.

The direct Hopper adapter's inverse scales match Transformer Engine exactly,
and reconstruction MSE equals
the separate RHT-then-quantize reference for FP16 and BF16. Fusion makes the
RHT-plus-quantization stage 1.27--1.34$\times$ faster and the complete forward
Linear pipeline 1.02--1.10$\times$ faster than the separate implementation.
For square-8192, Llama up, and Llama down, the fused rotation adds only 4.3\%,
6.5\%, and 2.4\% over block-FP8 Linear without rotation.

The bidirectional writer matches the two-writer implementation byte-for-byte
and scale-for-scale. Two-view preprocessing is 1.20--1.37$\times$ faster than
separate RHT plus TE quantization. Forward ranges from 1.01--1.12$\times$ and
forward plus backward from 0.99--1.04$\times$. Nsight Systems measures a
9.6\% reduction in the square-4096 GPU interval with unchanged GEMM and inverse
RHT costs.

For the block-level evaluation, we use a Llama-shaped MLP with a combined
$4096\rightarrow22016$ FC1, Transformer Engine fused SwiGLU, and a
$11008\rightarrow4096$ down projection. We pair every activation RHT with the
same transform on the corresponding weight input axis, using
$(xR)(wR)^\top=xw^\top$. On Hopper, directly folding SwiGLU into the 128-by-128
bidirectional writer is register-bound, so our final hybrid keeps TE SwiGLU in
forward and fuses inverse RHT with dSwiGLU in backward. Across three independent
processes, median speedups over our separate Triton RHT plus TE are 1.054$\times$ forward,
1.020$\times$ backward, and 1.027$\times$ for forward plus backward.
The respective process ranges are 1.053--1.059$\times$,
1.009--1.021$\times$, and 1.015--1.034$\times$. FP32 paired outputs and gradients agree within
$4.5\times10^{-7}$--$5.7\times10^{-7}$ relative L2, while BF16 differs by about
0.6\%. Paired TE block FP8 has 6.56--6.82\% relative error against BF16,
marginally below plain block FP8 for the output and all three gradients. We
report these as separate H100 NVL architecture results rather than mixing their
absolute latency with the H100 80GB HBM3 tables.

The preceding project-RHT comparison is an internal fusion ablation. Our
functionality-matched external baseline instead uses the unmodified Tri Dao H16
kernel at pinned commit \texttt{e7706faf}, a separate stock PyTorch BF16
multiply for the fixed Rademacher signs, and unmodified TE 2.18 block FP8. The
external and project RHT outputs match bit-for-bit in BF16.

\begin{table}[t]
\caption{BF16 Llama-shaped MLP latency on H100 NVL with 4096 tokens, hidden
size 4096, and intermediate size 11008. The external path composes unmodified
Tri Dao H16, stock PyTorch signs, and unmodified TE block FP8. Values are
medians across three independently started processes.}
\label{tab:rht-external}
\centering
\small
\begin{tabular}{lrrr}
\toprule
Phase & External RHT + TE (ms) & Fused (ms) & Speedup \\
\midrule
Forward & 4.533 & 1.964 & 2.308$\times$ \\
Backward & 6.171 & 3.488 & 1.769$\times$ \\
Forward + backward & 10.877 & 5.508 & 1.974$\times$ \\
\bottomrule
\end{tabular}
\end{table}

The forward, backward, and combined process ranges are 2.306--2.309$\times$,
1.767--1.775$\times$, and 1.969--1.981$\times$. Profiling attributes 5.119 ms
to four upstream H16 launches and 0.370 ms to four separate sign multiplies;
TE GEMM time is unchanged. Relative to TE block FP8 with RHT omitted, the fused
path is 0.941$\times$ in forward and 0.946$\times$ for forward plus backward.
The 1.974$\times$ result therefore measures an RHT-enabled pipeline, not a
speedup over omitting RHT.

\utkarsh{TODO: Repeat the RHT comparison across sign seeds or per-layer masks
and report confidence intervals rather than only process ranges.}

For dynamic weights, both quantized views and all valid scales match the
Transformer Engine reference exactly, while the
mapped gradient has $1.85\times10^{-9}$ relative L2 error and the AdamW update
is identical to a dense reference. In a controlled 120-step teacher--student
SwiGLU regression, dynamic RHT block FP8 ends at 0.001847 MSE versus 0.001907
for ordinary block FP8 and 0.000331 for BF16. This supports no additional
convergence regression from RHT relative to block FP8, but does not establish
BF16 parity. At the $4096\rightarrow11008$ shape, directly quantized dynamic
weights take 7.732 ms versus 7.814 ms for the interleaved BF16-working dynamic
control (1.011$\times$ faster). Nsight projects 6.106 versus 6.375 ms with 11
versus 13 GPU operations because two weight-quantization kernels are removed.
Both dynamic routes remain approximately 18--19\% slower than paired static
controls. These measurements exclude AdamW kernels and identify dynamic
materialization and gradient mapping as the remaining systems cost.

The fused optimizer passes an independently allocated two-path check: quantized
weight bytes, outputs, losses, and working gradients match exactly; after one
update, master and moment relative L2 errors are at
most $4.8\times10^{-8}$ and $2.4\times10^{-7}$. In the exact one-to-one
pipeline, forward is 1.000$\times$ and backward 0.997$\times$. Across three
processes for a 4096-token $4096\rightarrow22016\rightarrow4096$ SwiGLU MLP,
fusion reduces the median complete step from 11.233 to 8.356 ms
(1.370$\times$, range 1.344--1.399$\times$) and the optimizer tail from 4.224
to 1.338 ms (3.132$\times$) relative to explicit $R^\top$ mapping plus foreach
PyTorch AdamW. In a matched 120-step optimizer A/B, both paths begin at the same
loss; fused and mapped AdamW finish at 0.001817 and 0.001847 MSE, while median
step latency improves from 1.221 to 1.029 ms (1.186$\times$). This remains a
controlled regression rather than language-model pretraining.
Nsight Systems gives consistent architectural evidence: the complete profiled
range falls from 67 GPU operations and 9.525 ms projected GPU interval to 11
operations and 6.810 ms, while the optimizer-only range falls from 58
operations/3.903 ms to two operations/1.172 ms (one fused update per weight).

\utkarsh{TODO: Repeat convergence across seeds and add model-derived
activations plus language-model time-to-quality.}

Separately, we compare with ordinary Transformer Engine block FP8 without RHT,
using matched
original-basis FP32 masters and foreach AdamW. Across three independent
processes, median speedups are 0.842$\times$ for forward, 0.952$\times$ for
backward, and 0.898$\times$ for forward plus backward, showing that the RHT
itself remains overhead relative to an unrotated TE layer. The optimizer gain
reverses the complete result: the full step takes 7.859 ms versus 10.295 ms,
or 1.312$\times$ speedup (process range 1.310--1.318$\times$). Output relative
L2 against an independent BF16 computation is 0.0661 for RHT and 0.0662 for
plain TE. Thus the current end-to-end Hopper gain comes specifically from
making the randomized-transform optimizer path cheaper than ordinary TE's
weight quantization, gradient materialization, and AdamW sequence.
This is a system-level net-cost comparison, not the operation-matched optimizer
ablation above.

\utkarsh{TODO: Move the non-operation-matched 1.312x system comparison to the
appendix or replace it with an optimizer-matched off-the-shelf baseline.}

A shape sweep further shows that model width alone does not determine the
benefit. For a high-expansion $3584\rightarrow18944$ MLP, speedup grows from
1.020$\times$/1.008$\times$ forward/training at 1024 tokens to
1.073$\times$/1.032$\times$ at 8192 tokens. A $4096\rightarrow14336$ MLP at
4096 tokens reaches 1.059$\times$/1.021$\times$, whereas a larger
$8192\rightarrow28672$ case at 2048 tokens reaches only
1.025$\times$/1.012$\times$. This suggests that high FFN expansion and enough
tokens to saturate the fused writer are more favorable than hidden width alone.

\section{Limitations}
The primary hybrid MLP and final direct-TE comparison have three independent
process repetitions, while remaining timing sweeps still require confidence
intervals, and clock/power metadata. The H100 cannot provide native MXFP8 or NVFP4 throughput.
The patched Arthur baseline must be identified as a local BF16 compatibility
patch rather than unmodified upstream. The convergence experiment is a small
controlled regression, not language-model pretraining; model-scale quality and
time-to-accuracy remain to be established.

\utkarsh{TODO: Add native Blackwell MXFP8/NVFP4 throughput and model-scale
time-to-quality before submission.}

\section{Conclusion}
Tensor cores can implement a normalized FWHT at the H100 memory roofline while
reducing dynamic instructions. In a Llama-shaped MLP, fused RHT--FP8 is
1.974$\times$ faster than the functionality-matched off-the-shelf composition,
but adds 5.7\% latency relative to TE with RHT omitted. In a separate exact
optimizer ablation, fusing inverse RHT into AdamW reduces the complete step by
1.370$\times$.
The next question is whether this systems gain translates to model-scale
time-to-quality and survives native MXFP8/NVFP4 execution on Blackwell.

\bibliography{references}
\bibliographystyle{iclr2027_conference}

\section*{Reproducibility statement}
The supplementary artifact will include pinned baseline revisions, build
commands, correctness tests, raw JSON timing data, derived CSV tables, Nsight
Compute reports, and scripts that reproduce every table. Hardware and software
provenance are recorded alongside each result. The anonymous artifact will
exclude identifying paths and scheduler metadata.

\section*{Ethics statement}
This work studies GPU kernel efficiency and low-precision training. We do not
identify direct human-subject, privacy, or safety risks beyond the general
energy and hardware-resource impacts of large-scale machine-learning research.

\end{document}
